\section{法円作用の分岐係数と共有 annulus の分母剛性}
\label{mg8:section}

個々の core に作用座標を構成できても，その charts を大域 atlas に接続できるとは限らない．
法円作用の三枚の sheets における微分と，四 puncture neck の両端の annular periods を比較する．
これにより，混合五次では分母によらない大域障害が，六次では分母の剛性が得られる．
ここでは法円作用そのものが ambient-affine であるとは仮定しない．

\subsection{法円作用の分岐係数}

三価 branch の saturated star lattice を
\[
 u_1+u_2+u_3=\kappa v
\]
と表示する．$v$ は branch の primitive tangent，$u_a$ は各 sheet に向かう
primitive transverse lifts である．共通 Hamiltonian circle の作用関数 $h$ は，
各 sheet で整数アフィンであり，branch 上で連続である．
各 sheet におけるその線形部分を $\dd h_a$ とし，branch tangent での共通値を
$\dd h(v)$ と書く．次の整数を定める．
\begin{equation}\label{mg8:defect}
 \mathfrak d(h)=\sum_{a=1}^3\dd h_a(u_a)-\kappa\dd h(v).
\end{equation}
これは作用関数の加法定数を用いない．

% A short optional remark to place immediately after the definition of D(h).
\begin{remark}[Tropical divisor としての分岐係数]
\label{div9:interpretation}
係数 $\mathfrak d(h)$ は新しい一般的な tropical divisor の定義ではない．
三つの maximal cones の重みを一とし，その primitive normal vectors の lifts を
$u_a$ とすれば，Allermann--Rau の係数公式
\cite[Definition~3.4]{AllermannRau}は本稿の式に一致する．
lift の選択によらないことも同文献の Remark~3.5 に含まれる．
本稿で必要なのは，この既知の係数を normal circle の作用に適用し，
実際の SNC compactification の self-intersections と annular areas から
core ごとの総和を計算し，対向する normal circles の符号と比較することである．

各 core の頂点で用いる normal actions は異なる circle generators に由来する．
従って，それらの全係数の和を一つの大域関数の principal divisor の degree と
同一視することはできない．本稿の相殺は，branch ごとに二つの normal actions が
定数を除いて互いに負となることから導く．この議論は nodes を横断せず，
focus punctures を持つ底に対して compact tropical manifold の交点定理を
無条件に適用することもない．
\end{remark}


\begin{lemma}[分岐係数の不変性]\label{mg8:invariance}
整数 $\mathfrak d(h)$ は，$u_a$ の lift，$v$ の向き，sheet の順序によらない．
ambient integer affine transition によっても保存される．
従って，追加の branch singularities のない連結な開 branch に沿って一定である．
対向する二つの core の正の normal circles が同じ neck で互いに逆である場合には，
その両端で計算した係数は
\begin{equation}\label{mg8:end-cancellation}
 \mathfrak d_i(h_i)+\mathfrak d_j(h_j)=0
\end{equation}
を満たす．
\end{lemma}

\begin{proof}
$u_a$ を $u_a+m_av$ に替えると，$\kappa$ は $\kappa+\sum_a m_a$ に替わる．
\eqref{mg8:defect} の二項の変化は等しいので打ち消し合う．
$v$ の向きを逆にすれば $\kappa$ と $\dd h(v)$ の符号が同時に変わる．
sheet の置換にも明らかに不変である．

一つの ambient transition の線形部分を $T$ とする．
分岐方向を合わせ，transverse lifts を $T$ で運べば，
$\dd h_a(u_a)$ と relation の両方が保存される．
任意の lifts との差は前段により無関係である．
action cycle は各正則 sheet で平行であるから，その作用微分は整数アフィン座標で定数である．
従って係数は各 chart の中で一定であり，overlap の不変性から開 branch 全体で一定となる．

最後に，二つの core で取る正の normal circles は $\alpha_j=-\alpha_i$ である．
従って neck 上の二つの Hamiltonians は定数を除き $h_j=-h_i$ となる．
\eqref{mg8:defect} は $\dd h$ に線形なので，\eqref{mg8:end-cancellation} を得る．
特に，$T v_i=-v_j$，$T u_{i,a}=u_{j,\sigma(a)}+m_av_j$ と書く場合は，
$\sum_a m_a=-\kappa_i-\kappa_j$ を用いて同じ等式を直接検算できる．
\end{proof}

この係数が零であることは，共通 $h$ がその star で一つの ambient affine function に
延長するための条件である．しかし以下の障害は，各係数が零であることを要求せず，
対向する二端の相殺だけを使う．従って命題~\ref{quad:strong-obstruction} より強い比較になる．


% v9. General local denominator theorem, cone-index refinement,
% and all-plane sextic consequence. No new bibliography is required.
% Uses existing macros \ZZ and theorem environments.

\subsection{四 puncture neck の分母剛性}
\label{q9:four-puncture-section}

一つの core の annular equation には，高分母の整数解がある．
しかし，同じ actual neck の両端を比較すると，四 puncture 曲線には追加の整除条件が生じる．
以下では作用関数を微分して比較するので，primitive の変更による加法定数や，
有限の平滑化 parameter における作用切片の選択には依存しない．

\begin{proposition}[四 puncture neck の整除条件]
\label{q9:four-puncture-divisibility}
\label{q9:rigidity}
同じ marked actual neck の両端にある rational divisor components を
$C^-,C^+$ とし，それぞれの compactification における自己交点を $e^-,e^+\in\ZZ$ とする．
各曲線は四つの marked punctures を持ち，同じ $2+2$ 分割を実現する二つの
三価 saddle vertices $A,B$ と一つの内部 annulus からなる Reeb tree を用いる．
共有する二本の開 branch には追加の branch singularities を入れず，
その primitive action tangent lattices を接着するとする．
normal circles は同じ neck で互いに逆であり，puncture の正の meridian markings を保つとする．
両端には本稿の比例曲率 collars と Liouville cone construction を用い，
作用を $I_\gamma=(2\pi)^{-1}\int_\gamma\lambda$ と正規化する．

両端で作用座標に適合する rational cone charts が与えられ，各 branch lattice が三枚の
incident sheets の格子で飽和していると仮定する．
各内部 annulus の二本の primitive endpoint directions が，その regular torus の
action tangent lattice に生成する部分格子の index を $m^-,m^+\geq1$ とする．
このとき，二本の共有 branch の normal-action denominators は両端で一致し，
正の奇数 $q_A,q_B$ として
\begin{equation}\label{q9:cone-index-divisibility}
 q_Aq_B\mid m^-+m^+
\end{equation}
を満たす．特に，両方の annuli が一つずつの unimodular cone に対応する場合には
\begin{equation}\label{q9:four-puncture-q-one}
 q_A=q_B=1.
\end{equation}
\end{proposition}

\begin{proof}
共通の正の normal meridian を各端で $\alpha^\pm$ と書く．
各 core から外向きの primitive branch tangent $v_s^\pm$，$s\in\{A,B\}$，に対し
$\dd I_{\alpha^\pm}(v_s^\pm)=-q_s^\pm$ と置く．
共有 branch の整数アフィン同定は $v_s^-\mapsto-v_s^+$ であり，
$\alpha^+=-\alpha^-$ である．従って微分を比較すると $q_s^-=q_s^+$ となる．
その共通値を $q_s$ と書く．$A,B$ が交換される marking を用いた場合も，
片側の二頂点を先に交換すれば同じ記述になる．

各 puncture $\ell$ を囲む正の base meridian の degree-one lift を $\beta_\ell^\pm$ とし，
その puncture が接する saddle を $s(\ell)$ とする．primitive ray parameter $t$ で
\[
 I_{\alpha^\pm}=-q_{s(\ell)}t,
 \qquad I_{\beta_\ell^\pm}=-p_\ell^\pm t,
 \qquad p_\ell^\pm\in\ZZ
\]
と書ける．ここで core chart の加法定数は零に正規化してよく，
以下の比較ではこの表示の微分だけを使う．
三価 critical block の二つの arm cycles を $b_1,b_2$ とすれば，
第三の arm cycle は向きを揃えて $b_1+b_2$ である．
三つの sheets で branch direction が primitive となるためには，
その二つの numerators と和が全て $q_s$ と互いに素でなければならない．
$q_s$ が偶数なら二つの numerators は奇数で，和は偶数となり矛盾する．
従って $q_A,q_B$ は奇数である．

次に一つの端を固定し，記号 $\pm$ を省略する．
$A$ 側，$B$ 側の二つずつの puncture numerators の和を $P_A,P_B$ と置く．
puncture-compatible framing において，二組の puncture lifts の和を内部 annular torus
へ運ぶと，その実際の整数一次ホモロジーの中で
\[
 \beta_A+\beta_B=-e\alpha
\]
となる．これは normal circle bundle の Euler clutching による関係であり，
補集合の homology に移してから closed でない primitive を積分する操作ではない．
$(\alpha,\beta_A)$ は整数 torus basis である．
この基底で両 endpoint directions を並べた行列は
\[
 \begin{pmatrix}
 -q_A&-q_B\\
 -P_A&e q_B+P_B
 \end{pmatrix}.
\]
内部 annulus の面積が正であることと相対 Stokes の公式から
\[
 \frac{\operatorname{area}(E)}{2\pi w}
 =e+\frac{P_A}{q_A}+\frac{P_B}{q_B}>0
\]
を得る．従って上の determinant は負であり，その絶対値が $m$ であることは
\begin{equation}\label{q9:annular-index-equation}
 e q_Aq_B+q_BP_A+q_AP_B=m,
 \qquad
 \frac{\operatorname{area}(E)}{2\pi w}=\frac{m}{q_Aq_B}
\end{equation}
と同値である．normal Euler framing を整数だけ変更しても，この計算と以下の整除性は変わらない．

一方，共有 puncture の両側の cycles は，実際の regular torus の中で
\[
 \alpha^+=-\alpha^- ,\qquad
 \beta_\ell^+=\beta_\ell^-+k_\ell\alpha^-,
 \qquad k_\ell\in\ZZ
\]
と比較される．base puncture の向きが同じなので $\beta_\ell^-$ の係数は $+1$ である．
この式を共通 branch tangent に沿って微分すると
\begin{equation}\label{q9:puncture-integrality}
 \frac{p_\ell^-+p_\ell^+}{q_{s(\ell)}}=-k_\ell\in\ZZ
\end{equation}
となる．transport path を変更して node の周りの monodromy を加えても，
現在の meridian-node model では $k_\ell$ を整数だけ変更するので，整除性は変わらない．

両端で \eqref{q9:annular-index-equation} を $q_Aq_B$ で割り，足し合わせると
\[
 \frac{m^-+m^+}{q_Aq_B}
 =e^-+e^+
  +\sum_{\ell=1}^4\frac{p_\ell^-+p_\ell^+}{q_{s(\ell)}}
 =e^-+e^+-\sum_{\ell=1}^4 k_\ell\in\ZZ.
\]
これが \eqref{q9:cone-index-divisibility} である．
$m^-=m^+=1$ なら正の奇数 $q_Aq_B$ は二を割るので一であり，
\eqref{q9:four-puncture-q-one} を得る．
\end{proof}

\begin{remark}[Node 配分による表示]\label{q9:node-budget}
本稿の ordered resolution と puncture framing では，puncture basin $\ell$ に入る
meridian nodes の数を $n_\ell$ とすると $k_\ell=n_\ell-1$ である．
従って同じ必要条件は
\begin{equation}\label{q9:external-node-budget}
 \sum_{\ell=1}^4 n_\ell
 =4+e^-+e^+-\frac{m^-+m^+}{q_Aq_B}
\end{equation}
と書ける．mixed quintic の plane--quadric conic では $(e^-,e^+)=(4,2)$ なので，
unimodular の場合に右辺は $10-2/(q_Aq_B)$ である．
全平面 sextic の共有直線では $(e^-,e^+)=(1,1)$ なので $6-2/(q_Aq_B)$ となる．
前者の総 node 数は十，後者は六であり，分母剛性の後にはどちらも内部 annulus に二 nodes が残る．
この総数だけでは四つの external basins への個別の配分は決まらない．
命題の証明自体は $k_\ell$ の整数性だけを使い，node 数の非負性やこの具体的な framing 値を必要としない．
\end{remark}

\begin{corollary}[全平面 sextic の分岐分母]\label{q9:sextic-rigidity}
全平面六次退化の既存の cores と branch graph を保ち，
各 core に標準 Petersen fan の coarse cones への作用適合 charts を置く．
それらが大域作用適合 matroidal atlas に接続するとすれば，全三十本の branch の
normal-action denominator は一である．
これは既に構成した weights や node 配分に限定された結論ではない．
\end{corollary}

\begin{proof}
十五本の共有直線はそれぞれ四 punctures を持ち，二つの saddles と一つの内部 annulus を与える．
両端の Petersen chart で内部 annulus は一つの unimodular coarse cone に対応する．
命題~\ref{q9:four-puncture-divisibility} を各共有直線に適用すれば，
対応する二本の branch の分母がともに一となる．
\end{proof}

\begin{remark}[適用範囲と高分母に必要な変更]\label{q9:scope}
以上は，actual normal circles と puncture markings，二 saddle の Reeb tree，
共有 branch の primitive lattice，および各内部 annulus が一つの cone に対応する条件を使う．
各 sheet の作用微分の比較から従うため，各 normal Hamiltonian の ambient-linear 延長は仮定しない．
三 puncture 曲線には内部 annulus がなく，この剛性命題をそのまま適用することはできない．
高分母を得るために非 unimodular な内部 cones を許す場合でも，
\eqref{q9:cone-index-divisibility} により $m^-+m^+\geq q_Aq_B$ が必要である．
annulus の subdivision，Reeb tree の変更，追加の branch singularities を許す問題では，
その変更に応じて整合条件を立て直す必要がある．
\end{remark}

